\chapter*{Conclusion g\'en\'erale et perspectives}
\addcontentsline{toc}{chapter}{Conclusion g\'en\'erale et perspectives}
\markboth{Conclusion g\'en\'erale}{Conclusion g\'en\'erale}

% ---------------------------------------------------------------------
% GUIDE DE REDACTION
%   1. Rappeler en une phrase le contexte et l'objectif du projet.
%   2. Resumer le travail realise (sans entrer dans le detail technique).
%   3. Mettre en avant les apports / competences acquises (techniques
%      et personnelles).
%   4. Mentionner les limites actuelles.
%   5. Proposer des perspectives d'amelioration (cf. DOCUMENTATION_IA_NLP.md
%      section 13 pour des pistes deja identifiees sur le module NLP :
%      gestion ET/OU, synonymes douaniers, phrases negatives, regles
%      temporelles, historique des conversions, mode vocal, etc.)
% ---------------------------------------------------------------------

[\`A compl\'eter : conclusion g\'en\'erale rappelant le contexte, le travail
r\'ealis\'e (plateforme E-FORCE : gestion des processus BPMN, moteur de
r\`egles m\'etier avec g\'en\'eration DRL, module IA/NLP, gestion des
dossiers import/export), les comp\'etences acquises et les limites
actuelles.]

\textbf{Perspectives}

\begin{itemize}
    \item [\`A compl\'eter : gestion des connecteurs ET/OU dans le module
    NLP] ;
    \item [\`A compl\'eter : enrichissement du dictionnaire m\'etier
    (synonymes douaniers, terminologie UEMOA/CEMAC)] ;
    \item [\`A compl\'eter : historique des conversions NLP, mode vocal,
    support multilingue] ;
    \item [\`A compl\'eter : autres pistes d'\'evolution propres au
    projet].
\end{itemize}

\clearpage
